% 【本文件写什么】api-v0 契约层，叶子，只写语义不写实现
% - crate 路径：os/components/wateros-fs/fs-procfs/procfs-api/api-v0/src/
% - 列出并说明：pub trait、核心类型、错误枚举、常量
% - 每个公开 API 的前置条件、后置条件、线程/中断上下文限制
% - 与 Linux/用户态约定的对齐点与 deliberate 差异
% - v0 版本当前行为与后续可替换 extension 点
% 【事实来源】
% - 上述 crate 下全部 .rs 源文件
% - docs/exports/public-api/ 中对应符号表，以源码为准
% 【不写什么】
% - impl 内部数据结构、汇编、具体算法步骤

\subsubsection{api-v0}

\paragraph{类型与回调。}
\code{TaskId}与\code{task::TaskId}数值一致，api层不依赖\code{wateros-task}。\code{ProcMountLine}描述\code{/proc/mounts}单行，含 \code{device}、\code{mount\_point}、\code{fstype}。三类函数指针回调：
\begin{itemize}
 \item \code{TaskArgvLookup}：按leader task id查询\code{argv}；
 \item \code{TaskExeLookup}：查询\code{exe}路径；
 \item \code{MountListLookup}：枚举当前挂载表行。
\end{itemize}
由VFS层在\code{mount\_procfs\_at}前注册，impl-kernel通过静态\code{Mutex<Option<...>>}保存。

\paragraph{\code{ProcFsView}契约。}
\code{rel\_path}为相对\code{/proc}的路径，可带或不带前导 \code{/}。方法语义：
\begin{itemize}
 \item \code{exists}：已知proc节点，含目录与文件返回\code{true}；
 \item \code{metadata}：返回\code{FsMetadata}快照；不存在为\code{NotFound}；
 \item \code{read}：读普通文件；目录路径返回\code{NotAFile}；
 \item \code{read\_dir}：列目录；非目录返回\code{NotAFile}。
\end{itemize}
错误类型复用\code{FsError}。字段内容与完整Linux\code{/proc}有意保持LTP/bring-up子集。

\paragraph{扩展点。}
v0通过回调注入task与挂载信息，避免procfs api直接依赖VFS或task crate；后续可增补节点类型而不改\code{ProcFsView}方法表，由impl内路径解析扩展。
